\documentclass[12pt,a4paper]{article}

%% ── Packages ─────────────────────────────────────────────────────────────────
\usepackage[a4paper, top=2.5cm, bottom=2.5cm, left=2.5cm, right=2.5cm]{geometry}
\usepackage{graphicx, float, booktabs, longtable, hyperref, xcolor}
\usepackage{enumitem, fancyhdr, titlesec, array, amsmath, microtype}
\usepackage[T1]{fontenc}
\usepackage[utf8]{inputenc}
\usepackage{lmodern}

%% ── Color palette ────────────────────────────────────────────────────────────
\definecolor{sigmaBlue}{RGB}{21,101,192}
\definecolor{sigmaGreen}{RGB}{46,125,50}
\definecolor{sigmaRed}{RGB}{198,40,40}
\definecolor{sigmaOrange}{RGB}{230,81,0}
\definecolor{sigmaGray}{RGB}{117,117,117}
\definecolor{abstractBg}{RGB}{240,244,255}

%% ── Hyperlinks ───────────────────────────────────────────────────────────────
\hypersetup{colorlinks=true, linkcolor=sigmaBlue, urlcolor=sigmaBlue,
  pdftitle={SigmaFlow --- Automated DMAIC Analysis Report}, pdfauthor={SigmaFlow Engine v8}}

%% ── Section style ────────────────────────────────────────────────────────────
\titleformat{\section}{\color{sigmaBlue}\Large\bfseries}{\thesection}{1em}{}
\titleformat{\subsection}{\color{sigmaBlue}\large\bfseries}{\thesubsection}{1em}{}
\titleformat{\subsubsection}{\color{sigmaGray}\normalsize\bfseries}{\thesubsubsection}{1em}{}

%% ── Header / Footer ──────────────────────────────────────────────────────────
\pagestyle{fancy}
\fancyhf{}
\lhead{\small\color{sigmaGray}SigmaFlow v8 — Statistical Analysis Report}
\rhead{\small\color{sigmaGray}09 March 2026, 21:41}
\cfoot{\small\thepage}

%% ── Title ────────────────────────────────────────────────────────────────────
\title{
  \vspace{1.5cm}
  {\color{sigmaBlue}\Huge\bfseries SigmaFlow --- Automated DMAIC Analysis Report}\\[0.6em]
  {\large\color{sigmaGray} Automated Lean Six Sigma / DMAIC Analysis}\\[2em]
  {\normalsize\textbf{Engine:} SigmaFlow Engine v8}
}
\author{}
\date{09 March 2026, 21:41}

\begin{document}
\maketitle
\thispagestyle{empty}
\newpage

%% ══════════════════════════════════════════════════════════════════════════════
\section*{Abstract}
\addcontentsline{toc}{section}{Abstract}
%% ══════════════════════════════════════════════════════════════════════════════

\begin{flushleft}
\colorbox{abstractBg}{\parbox{0.96\textwidth}{
\smallskip
Automated SPC analysis of 'manufacturing\_process\_dataset' (300 observations, 7 variables) by SigmaFlow v10. 1 statistical anomaly(ies) detected (1 critical, 0 warnings). Regression R?=0.084. 1 column(s) flagged as non-normal.
\smallskip
}}
\end{flushleft}

\newpage
\tableofcontents
\newpage

%% ══════════════════════════════════════════════════════════════════════════════
\section{Introduction}
%% ══════════════════════════════════════════════════════════════════════════════


This report was automatically generated by \textbf{SigmaFlow v8}, a modular
Python framework for industrial process analysis. A total of
\textbf{1 dataset(s)} were processed through the complete analysis pipeline.

\medskip

The analysis pipeline consists of the following stages:

\begin{enumerate}[leftmargin=2em]
  \item \textbf{Dataset Detection} — Automatic classification of dataset type
        (SPC, capability, logistics, Pareto, service, DOE).
  \item \textbf{Statistical Analysis} — Descriptive statistics, control limits,
        capability indices (Cp, Cpk), and trend detection.
  \item \textbf{Visualization} — Control charts, histograms, Pareto charts,
        capability plots, and correlation heatmaps.
  \item \textbf{Root Cause Detection} — Pearson and Spearman correlation analysis
        to identify variables most associated with defects or variation.
  \item \textbf{Insight Generation} — Western Electric Rules and statistical
        rules evaluated against the data; insights exported as JSON.
  \item \textbf{Report Generation} — This document: fully annotated with
        interpretations and recommendations.
\end{enumerate}

\begin{center}
\begin{tabular}{lllcc}
\toprule
\textbf{Dataset} & \textbf{Type} & \textbf{Shape} & \textbf{Time (s)} & \textbf{Insights} \\
\midrule
  manufacturing\_process\_dataset & SPC & 300\(\times\)7 & 8.38s & 2\\
\bottomrule
\end{tabular}
\end{center}

\newpage

%% ══════════════════════════════════════════════════════════════════════════════
\section{Methodology}
%% ══════════════════════════════════════════════════════════════════════════════


The following statistical methods and tools are employed in this analysis:

\subsubsection*{Statistical Process Control (SPC)}

Control charts are used to distinguish between \textit{common-cause variation}
(inherent randomness) and \textit{special-cause variation} (assignable, non-random
events). The \textbf{XmR chart} (Individuals and Moving Range) is used for
individual observations. Control limits are set at $\mu \pm 3\sigma$.

\subsubsection*{Western Electric Rules}

Four detection rules are applied to identify non-random patterns:

\begin{itemize}[leftmargin=2em]
  \item \textbf{Rule 1} — Any single point outside $\pm 3\sigma$ limits
        (probability $\approx$ 0.27\% under normality).
  \item \textbf{Rule 2} — Nine consecutive points on the same side of the
        center line (probability $\approx$ 0.4\%).
  \item \textbf{Rule 3} — Six consecutive points trending upward or downward
        (indicates drift).
  \item \textbf{Rule 4} — Fourteen consecutive points alternating direction
        (indicates over-adjustment or dual-stream process).
\end{itemize}

\subsubsection*{Process Capability Indices}

When specification limits are available, the following indices are computed:

\begin{center}
\begin{tabular}{lll}
\toprule
\textbf{Index} & \textbf{Formula} & \textbf{Interpretation} \\
\midrule
$C_p$  & $(USL - LSL) / 6\sigma$ & Potential capability (centered) \\
$C_{pk}$ & $\min(C_{pu}, C_{pl})$ & Actual capability (off-centering penalty) \\
$C_{pu}$ & $(USL - \mu) / 3\sigma$ & Upper one-sided capability \\
$C_{pl}$ & $(\mu - LSL) / 3\sigma$ & Lower one-sided capability \\
\bottomrule
\end{tabular}
\end{center}

\medskip

\textbf{Capability thresholds:}
\begin{itemize}[leftmargin=2em]
  \item $C_{pk} < 1.00$ \textrightarrow\ Process \textcolor{sigmaRed}{\textbf{not capable}}
  \item $1.00 \le C_{pk} < 1.33$ \textrightarrow\ \textcolor{sigmaOrange}{\textbf{Marginal}} capability
  \item $1.33 \le C_{pk} < 1.67$ \textrightarrow\ \textcolor{sigmaGreen}{\textbf{Acceptable}} capability
  \item $C_{pk} \ge 1.67$ \textrightarrow\ \textcolor{sigmaGreen}{\textbf{Excellent}} (Six Sigma level)
\end{itemize}

\subsubsection*{Root Cause Analysis}

Pearson and Spearman rank correlations are computed between all process
variables and the primary quality output. Variables are ranked by absolute
correlation to identify the most likely process drivers.

\begin{center}
\begin{tabular}{ll}
\toprule
\textbf{|r| threshold} & \textbf{Strength} \\
\midrule
$|r| \ge 0.70$ & Strong association — high priority investigation \\
$0.50 \le |r| < 0.70$ & Moderate association — medium priority \\
$0.30 \le |r| < 0.50$ & Weak association — monitor \\
$|r| < 0.30$ & Negligible \\
\bottomrule
\end{tabular}
\end{center}

\newpage

%% ══════════════════════════════════════════════════════════════════════════════
\section{Dataset Analysis}
%% ══════════════════════════════════════════════════════════════════════════════

\subsection{manufacturing\_process\_dataset — SPC}
\subsubsection*{Dataset Overview}
\begin{description}[leftmargin=3.5cm, style=nextline]
  \item[Dataset]    \textbf{manufacturing\_process\_dataset}
  \item[Type]       \textbf{SPC}
  \item[Dimensions] 300 rows \(\times\) 7 columns
  \item[Processing] 8.38 seconds
\end{description}
\subsubsection*{Summary}
Automated SPC analysis of 'manufacturing\_process\_dataset' (300 observations, 7 variables) by SigmaFlow v10. 1 statistical anomaly(ies) detected (1 critical, 0 warnings). Regression R?=0.084. 1 column(s) flagged as non-normal.
\subsubsection*{Statistical Analysis Results}
\begin{longtable}{>{\bfseries}p{6.5cm} p{8.5cm}}
\toprule \textbf{Metric} & \textbf{Value} \\ \midrule
\endfirsthead
\multicolumn{2}{c}{\small(continued)} \\
\toprule \textbf{Metric} & \textbf{Value} \\ \midrule
\endhead
  n & 300 \\
  mean & 79.8672 \\
  std & 3.0066 \\
  x\_chart.UCL & 88.887 \\
  x\_chart.CL & 79.867 \\
  x\_chart.LCL & 70.848 \\
  x\_chart.n\_ooc & 2 \\
  mr\_chart.UCL & 11.26 \\
  mr\_chart.CL & 3.447 \\
  mr\_chart.n\_ooc & 2 \\
  trend.direction & stable \\
  trend.tau & -0.0175 \\
  trend.p\_value & 0.7629 \\
\bottomrule
\end{longtable}
\subsubsection*{Process Stability Analysis}
\begin{figure}[H]
  \centering
  \includegraphics[width=0.95\linewidth]{/home/osvaldinojunior/Documents/Projetos/SigmaFlow/output/figures/manufacturing_process_dataset/spc_xmr_chart.png}
  \caption{Spc Xmr Chart}
\end{figure}
\medskip
\begin{figure}[H]
  \centering
  \includegraphics[width=0.95\linewidth]{/home/osvaldinojunior/Documents/Projetos/SigmaFlow/output/figures/manufacturing_process_dataset/spc_trend.png}
  \caption{Spc Trend}
\end{figure}
\medskip
\subsubsection*{Graphical Results}
\begin{figure}[H]
  \centering
  \includegraphics[width=0.95\linewidth]{/home/osvaldinojunior/Documents/Projetos/SigmaFlow/output/figures/manufacturing_process_dataset/regression_coefficients.png}
  \caption{Regression Coefficients}
\end{figure}
\medskip
\begin{figure}[H]
  \centering
  \includegraphics[width=0.95\linewidth]{/home/osvaldinojunior/Documents/Projetos/SigmaFlow/output/figures/manufacturing_process_dataset/regression_actual_vs_predicted.png}
  \caption{Regression Actual Vs Predicted}
\end{figure}
\medskip
\begin{figure}[H]
  \centering
  \includegraphics[width=0.95\linewidth]{/home/osvaldinojunior/Documents/Projetos/SigmaFlow/output/figures/manufacturing_process_dataset/regression_diagnostics.png}
  \caption{Regression Diagnostics}
\end{figure}
\medskip
\subsubsection*{Detected Anomalies \& Statistical Insights}
\paragraph{\textcolor{sigmaRed}{Critical Findings}}
\textbf{Western Electric Rule 1}
\begin{itemize}[leftmargin=2em, itemsep=0pt]
  \item \textbf{Finding:} Rule 1: 2 point(s) outside the 3\(\sigma\) control limits (at indices [128, 174]).
  \item \textbf{Interpretation:} A data point beyond \(\pm\)3\(\sigma\) from the mean is statistically rare (\textasciitilde{}0.27\% chance under normality). Its presence strongly suggests a special-cause event --- something changed in the process at that moment.
  \item \textbf{Recommended Investigation:} Identify and investigate the specific points flagged. Check for equipment malfunction, operator changes, raw material variation, or measurement errors. Remove confirmed special causes before computing capability indices.
\end{itemize}
\medskip
\paragraph{\textcolor{sigmaGreen}{Informational}}
\textbf{Trend Stable}
\begin{itemize}[leftmargin=2em, itemsep=0pt]
  \item \textbf{Finding:} No statistically significant trend detected in the process data.
  \item \textbf{Interpretation:} The Spearman correlation coefficient (\(\tau\) = -0.018) and p-value (p = 0.7629) indicate no systematic drift in the process over time.
  \item \textbf{Recommended Investigation:} Continue routine monitoring. The process is temporally stable.
\end{itemize}
\medskip
\subsubsection*{Root Cause Indicators}
Correlation analysis was performed using the variable \textbf{defect} as the quality target.
\medskip
No variables showed meaningful correlation with 'defect'. The variation may be driven by factors not present in this dataset.
\paragraph{Variable Importance Ranking}
\begin{center}
\begin{tabular}{lrrl}
\toprule \textbf{Variable} & \textbf{Pearson r} & \textbf{Spearman r} & \textbf{Strength} \\
\midrule
  temperature\_C & \textcolor{sigmaGray}{\textbf{+0.263}} & +0.217 & Negligible \\\\
  pressure\_bar & \textcolor{sigmaGray}{\textbf{-0.079}} & -0.089 & Negligible \\\\
  machine\_speed\_rpm & \textcolor{sigmaGray}{\textbf{+0.058}} & +0.061 & Negligible \\\\
  dimension\_mm & \textcolor{sigmaGray}{\textbf{+0.020}} & -0.004 & Negligible \\\\
\bottomrule
\end{tabular}
\end{center}
\begin{figure}[H]
  \centering
  \includegraphics[width=0.95\linewidth]{/home/osvaldinojunior/Documents/Projetos/SigmaFlow/output/figures/manufacturing_process_dataset/correlation_heatmap.png}
  \caption{Correlation Heatmap}
\end{figure}
\medskip
\begin{figure}[H]
  \centering
  \includegraphics[width=0.95\linewidth]{/home/osvaldinojunior/Documents/Projetos/SigmaFlow/output/figures/manufacturing_process_dataset/variable_importance.png}
  \caption{Variable Importance}
\end{figure}
\medskip
\newpage

\newpage

%% ══════════════════════════════════════════════════════════════════════════════
\section{Recommendations}
%% ══════════════════════════════════════════════════════════════════════════════


Based on the automated analysis, the following actions are recommended:

\begin{enumerate}[leftmargin=2em]
  \item \textcolor{sigmaRed}{\textbf{Address 1 critical finding(s) immediately.}}
        Investigate all points flagged by Western Electric rules outside $3\sigma$ limits.
        Document findings and corrective actions.
  \item \textcolor{sigmaOrange}{\textbf{Review 0 warning(s)}} with the
        process engineering team. Schedule corrective actions within the next sprint cycle.
  \item \textbf{Re-run SigmaFlow} after implementing improvements to verify effectiveness
        of corrective actions.
  \item \textbf{Expand data collection} on variables identified in root cause analysis
        to confirm or refute correlation findings.
  \item \textbf{Set up real-time monitoring} for datasets with out-of-control processes.
\end{enumerate}



\medskip\noindent
\textit{All findings should be reviewed by a qualified process engineer before
implementing corrective actions. Correlation does not imply causation.}

\newpage

%% ══════════════════════════════════════════════════════════════════════════════
\section{Conclusion}
%% ══════════════════════════════════════════════════════════════════════════════


The SigmaFlow v8 automated analysis successfully processed
\textbf{1 dataset(s)} and produced the following summary:

\begin{center}
\begin{tabular}{lc}
\toprule
\textbf{Finding Category} & \textbf{Count} \\
\midrule
\textcolor{sigmaRed}{Critical} (immediate action) & 1 \\
\textcolor{sigmaOrange}{Warning} (monitor / address) & 0 \\
\textcolor{sigmaGreen}{Informational} (stable) & 1 \\
\bottomrule
\end{tabular}
\end{center}

\medskip

This report was generated automatically by \textbf{SigmaFlow v8}.
All outputs — figures, insights.json, report.tex, and report.pdf —
are available in the \texttt{output/} directory.
Statistical findings and recommendations should be reviewed by a
qualified Six Sigma practitioner or process engineer.

\end{document}